%! Equations and Inequalities — Unit 1, Lesson 1.2 — Warm-Up
\documentclass[10pt]{article}
\usepackage{algebra2-article}
\usepackage{algebra2-boxes}
\newcommand{\TallMath}[1]{$\displaystyle #1\rule[-1.4em]{0pt}{3.2em}$}

% Pre-drawn number line; #1 receives the shading/endpoint overlay.
\newcommand{\numline}[1]{%
  \begin{tikzpicture}[baseline=-3pt, x=0.58cm, y=0.58cm]
    \draw[<->, thick, charcoal] (-5.6,0) -- (5.6,0);
    \foreach \t in {-5,-4,-3,-2,-1,0,1,2,3,4,5} {\draw[charcoal] (\t,0.13) -- (\t,-0.13);}
    \foreach \t in {-4,-2,0,2,4} {\node[below=0pt, font=\scriptsize] at (\t,-0.10) {$\t$};}
    #1
  \end{tikzpicture}}

\begin{document}

\pageheader{Unit 1, Lesson 1.2}{Warm-Up}
\namedateperiod

\vspace{0.06in}

\begin{notesbox}{Getting Ready: Skills We'll Use Today}
\small Three skills, one from each of the last two lessons. Work quickly.

\vspace{3pt}
\textbf{1. Name the move (Lesson 1.0).} The work is done for you. Write the \emph{property} that
makes each line legal.
\begin{center}
\renewcommand{\arraystretch}{1.35}
\begin{tabularx}{\linewidth}{@{}l X@{}}
\hline
\textbf{Line} & \textbf{Why is this allowed?} \\
\hline
$5(x-3)=2x+9$ & \emph{given} \\
$5x-15=2x+9$ & \blank{7.6cm} \\
$3x-15=9$ & \blank{7.6cm} \\
$3x=24$ & \blank{7.6cm} \\
$x=8$ & \blank{7.6cm} \\
\hline
\end{tabularx}
\end{center}
``I moved it to the other side'' is not one of the answers. Every line has a \emph{name}.

\vspace{3pt}
\textbf{2. Factor completely (Lesson 1.1).}
\begin{center}
\renewcommand{\arraystretch}{1.5}
\begin{tabular}{@{}l l@{}}
(a) \; $x^2-5x-14=$ \blank{3.6cm} &
(b) \; $3x^2-27=$ \blank{3.6cm} \\
\end{tabular}
\end{center}
Now finish this sentence, which you already believe: \\
If two numbers multiply to give \textbf{zero}, then \blank{4.2cm} or \blank{4.2cm} must be zero.
\\[2pt]
Does the same reasoning work if the two numbers multiply to give \textbf{eight}? \blank{1.6cm}

\vspace{4pt}
\textbf{3. Read the picture.} Each number line below shows a \emph{solution set}. Write the
inequality it represents.
\begin{center}
\begin{tabular}{@{}c c@{}}
(a) & (b) \\[2pt]
\numline{\draw[ultra thick, forest] (-1,0) -- (5.55,0);
         \draw[thick, forest, fill=white] (-1,0) circle (2.7pt);}
&
\numline{\draw[ultra thick, forest] (3,0) -- (-5.55,0);
         \fill[forest] (3,0) circle (2.7pt);}
\\[3pt]
$x$ \blank{2.6cm} & $x$ \blank{2.6cm} \\
\end{tabular}
\end{center}
Roughly how many numbers are in solution set (a)? \blank{4.0cm} \\
That answer is exactly why we \emph{draw} a solution set instead of listing it.

\end{notesbox}

\end{document}
